% GENERATED FILE. Edit canonical lessons, then run npm run generate:books.
% canonical-source-hash: fnv1a64:9a2bbcd782a88a00
% canonical-lessons: TE-C292-gudise, TE-C292-kaluva, TE-C292-bhavanam, TE-C292-kota, TE-C292-karkhana

\chapter{Hut, Canal, Building, Fort, Factory}
\label{ch:te-hut-canal-building-fort-factory}
\hlchaptermodality{292}

\begin{chapteropening}
\textbf{By the end of this chapter:} \emph{I can say a hut, a canal, a building, a fort and a factory.}

Complete the last lesson of chapter 292: I can say a hut, a canal, a building, a fort and a factory.
\end{chapteropening}

\section[\texorpdfstring{\te{గుడిసె} (guḍise)}{గుడిసె (guḍise)}]{\te{గుడిసె} (guḍise) — a hut}

\label{lesson:TE-C292-gudise}



\begin{quote}
Before the new one: say the Telugu for opposite, then the Telugu for around.
\end{quote}

\subsection*{You'll want to know: \te{గుడిసె}}
\textbf{\te{గుడిసె}} — \emph{guḍise} — \textquotedblleft{}a hut\textquotedblright{}.

\begin{grammarlens}[title={What you've built}]
One more word for this chapter.
\end{grammarlens}

\subsection*{Your turn}
\begin{itemize}
\raggedright
  \item \emph{Say it:} \emph{guḍise}
  \item \emph{Say it:} \emph{guḍise}, once more
  \item \emph{Say it:} say \emph{guḍise}
  \item \emph{Recall:} say the Telugu for opposite, then the Telugu for around, then say \emph{guḍise} again
\end{itemize}

\begin{tcolorbox}[breakable,colback=teal!4,colframe=teal!35!black,title={Before you move on}]
What does \te{గుడిసె} mean? (\textquotedblleft{}A hut\textquotedblright{}.) Say it once more.
\end{tcolorbox}
\section[\texorpdfstring{\te{కాలువ} (kāluva)}{కాలువ (kāluva)}]{\te{కాలువ} (kāluva) — a canal}

\label{lesson:TE-C292-kaluva}



\begin{quote}
Before the new one: say the Telugu for around, then the Telugu for a hut.
\end{quote}

\subsection*{You'll want to know: \te{కాలువ}}
\textbf{\te{కాలువ}} — \emph{kāluva} — \textquotedblleft{}a canal\textquotedblright{}.

\begin{grammarlens}[title={What you've built}]
One more word for this chapter.
\end{grammarlens}

\subsection*{Your turn}
\begin{itemize}
\raggedright
  \item \emph{Say it:} \emph{kāluva}
  \item \emph{Say it:} \emph{kāluva}, once more
  \item \emph{Say it:} say \emph{kāluva}
  \item \emph{Recall:} say the Telugu for around, then the Telugu for a hut, then say \emph{kāluva} again
\end{itemize}

\begin{tcolorbox}[breakable,colback=teal!4,colframe=teal!35!black,title={Before you move on}]
What does \te{కాలువ} mean? (\textquotedblleft{}A canal\textquotedblright{}.) Say it once more.
\end{tcolorbox}
\section[\texorpdfstring{\te{భవనం} (bhavanaṁ)}{భవనం (bhavanaṁ)}]{\te{భవనం} (bhavanaṁ) — a building}

\label{lesson:TE-C292-bhavanam}



\begin{quote}
Before the new one: say the Telugu for a hut, then the Telugu for a canal.
\end{quote}

\subsection*{You'll want to know: \te{భవనం}}
\textbf{\te{భవనం}} — \emph{bhavanaṁ} — \textquotedblleft{}a building\textquotedblright{}.

\begin{grammarlens}[title={What you've built}]
One more word for this chapter.
\end{grammarlens}

\subsection*{Your turn}
\begin{itemize}
\raggedright
  \item \emph{Say it:} \emph{bhavanaṁ}
  \item \emph{Say it:} \emph{bhavanaṁ}, once more
  \item \emph{Say it:} say \emph{bhavanaṁ}
  \item \emph{Recall:} say the Telugu for a hut, then the Telugu for a canal, then say \emph{bhavanaṁ} again
\end{itemize}

\begin{tcolorbox}[breakable,colback=teal!4,colframe=teal!35!black,title={Before you move on}]
What does \te{భవనం} mean? (\textquotedblleft{}A building\textquotedblright{}.) Say it once more.
\end{tcolorbox}
\section[\texorpdfstring{\te{కోట} (kōṭa)}{కోట (kōṭa)}]{\te{కోట} (kōṭa) — a fort}

\label{lesson:TE-C292-kota}



\begin{quote}
Before the new one: say the Telugu for a canal, then the Telugu for a building.
\end{quote}

\subsection*{You'll want to know: \te{కోట}}
\textbf{\te{కోట}} — \emph{kōṭa} — \textquotedblleft{}a fort\textquotedblright{}.

\begin{grammarlens}[title={What you've built}]
One more word for this chapter.
\end{grammarlens}

\subsection*{Your turn}
\begin{itemize}
\raggedright
  \item \emph{Say it:} \emph{kōṭa}
  \item \emph{Say it:} \emph{kōṭa}, once more
  \item \emph{Say it:} say \emph{kōṭa}
  \item \emph{Recall:} say the Telugu for a canal, then the Telugu for a building, then say \emph{kōṭa} again
\end{itemize}

\begin{tcolorbox}[breakable,colback=teal!4,colframe=teal!35!black,title={Before you move on}]
What does \te{కోట} mean? (\textquotedblleft{}A fort\textquotedblright{}.) Say it once more.
\end{tcolorbox}
\section[\texorpdfstring{\te{కార్ఖానా} (kārkhānā)}{కార్ఖానా (kārkhānā)}]{\te{కార్ఖానా} (kārkhānā) — a factory}

\label{lesson:TE-C292-karkhana}



\begin{quote}
Before the new one: say the Telugu for a building, then the Telugu for a fort.
\end{quote}

\subsection*{You'll want to know: \te{కార్ఖానా}}
\textbf{\te{కార్ఖానా}} — \emph{kārkhānā} — \textquotedblleft{}a factory\textquotedblright{}.

\begin{grammarlens}[title={What you've built}]
One more word for this chapter.
\end{grammarlens}

\subsection*{Your turn}
\begin{itemize}
\raggedright
  \item \emph{Say it:} \emph{kārkhānā}
  \item \emph{Say it:} \emph{kārkhānā}, once more
  \item \emph{Say it:} say \emph{kārkhānā}
  \item \emph{Recall:} say the Telugu for a building, then the Telugu for a fort, then say \emph{kārkhānā} again
\end{itemize}

\begin{tcolorbox}[breakable,colback=teal!4,colframe=teal!35!black,title={Before you move on}]
What does \te{కార్ఖానా} mean? (\textquotedblleft{}A factory\textquotedblright{}.) Say it once more.
\end{tcolorbox}
